\originalchapter{超精细结构与外场}
\originalsection{角动量耦合与Zeeman区间}
核自旋$\mathbf I$和电子总角动量$\mathbf J$的磁偶极超精细作用写为$H_{\rm hf}=A\mathbf I\cdot\mathbf J/\hbar^2$. 弱场用$\mathbf F=\mathbf I+\mathbf J$, 得
\[E_F=\frac A2[F(F+1)-I(I+1)-J(J+1)].\]
$I,J$均至少为1时可另有电四极作用; $J=1/2$时没有该电子四极张量. Zeeman项写成$\mu_BB(g_JJ_z+g_II_z)/\hbar$, 这里$g_I$已用Bohr磁子作单位并包含核磁矩的实际符号.

Landé因子可由$\mathbf J$方向上的磁矩投影得到:
\[g_J=g_L\frac{j(j+1)+l(l+1)-s(s+1)}{2j(j+1)}+g_S\frac{j(j+1)-l(l+1)+s(s+1)}{2j(j+1)}.\]
把$J,I$再次投影到$F$上即得$g_F$, 公式结构相同. 强场使Zeeman项超过超精细耦合, 好量子数改为$m_J,m_I$; 此时$F$不再守恒.

\originalsection{Breit--Rabi公式的二乘二来源}
对$J=1/2$, 固定$m=m_F$的基为$\ket{m_I=m-1/2,m_J=1/2}$和$\ket{m_I=m+1/2,m_J=-1/2}$. 超精细对角元为$A(m/2-1/4)$与$A(-m/2-1/4)$; 升降算符给非对角元$A\sqrt{(I+1/2)^2-m^2}/2$. 加上Zeeman项后直接求两个根, 记$\Delta E=A(I+1/2)$及$x=(g_J-g_I)\mu_BB/\Delta E$, 得
\begin{equation}E_m^{\pm}=-\frac A4+g_I\mu_BmB\pm\frac{\Delta E}{2}\sqrt{1+\frac{2m}{I+1/2}x+x^2}.\label{eq:BR}\end{equation}
最外侧$m=\pm(I+1/2)$仅有一个基矢, 必须只取实际存在的分支. $x\to0$恢复$F=I\pm1/2$; $|x|\to\infty$恢复$m_I,m_J$. 在中间区间, 有些跃迁的两个一阶斜率相等, 形成对磁场噪声较不敏感的钟跃迁.

\originalsection{Stark作用、宇称与分子偶极}
电偶极作用为$H_E=-\mathbf d\cdot\mathbf E$. 宇称把$\mathbf d$反号, 所以非简并宇称本征态有$\langle\mathbf d\rangle=0$. 弱场一阶位移为零, 二阶为
\[\Delta E_a=-\frac12\alpha_a(0)E^2,\qquad\alpha_a(0)=2\sum_{b\ne a}\frac{|\bra b d_z\ket a|^2}{E_b-E_a}.\]
对基态分母为正, 极化率为正. 对激发态不能先验认定正号. 若两个异宇称能级相距$\hbar\omega_s$, 其矩阵为$\bigl(\begin{smallmatrix}\hbar\omega_s/2&-dE\\-dE&-\hbar\omega_s/2\end{smallmatrix}\bigr)$, 根为$\pm\sqrt{(\hbar\omega_s/2)^2+d^2E^2}$. 弱场平方, 强混合时近似线性; 不需要违反宇称守恒来解释.

分子的“永久偶极”通常指体轴上的电荷分离$d_0$. 自由转子在实验室中没有固定方向, 故定宇称定转动态的实验室偶极平均仍为零. 外场混合相反宇称的转动态并定向分子, 当$d_0E$超过转动间隔而仍远小于电子激发尺度时, 可以观察到近似线性Stark斜率. 磁矩是轴向矢量, 在空间反演下不变; 原子可以有磁偶极而不必有电偶极. 永久原子电偶极还涉及时间反演对称性, 与分子的体轴偶极不能混为一谈.
